\section{Unified synthetic comparison of frequency--relay screens}
\label{sec:p17-benchmark}
To distinguish screening conservatism from physical feasibility, we evaluated four methods on an identical, seeded set of \(90\) \emph{synthetic} piecewise-current scenarios (seed \(1717\)). Thirty scenarios each used full, hold and linear timer reset. Every scenario specified primary and counterfactual backup current templates, pickup and inverse-time multipliers, independent primary/backup current scales, breaker-delay intervals, finite clearing deadlines, and a coordination margin. A stipulated sign-aware frequency gate accepted only if \(f_0-\Delta f\ge f_{\min}\), with frequencies expressed in hertz.

The four methods were nominal-current evaluation; common-derating endpoint evaluation; the analytically exact independent-rectangle criterion under the fixed-template monotone-timer assumptions of Section~\ref{sec:p13-continuum}; and explicit time-domain evaluation on a finite grid of primary/backup derating and delay values. All were evaluated against the \emph{same} independent rectangular model reference. The resulting approval counts were \(41,25,22,22\), respectively; false approvals were \(19,3,0,0\), and conservative rejections were zero for all four methods. Agreement of the finite time-domain grid with the theorem-backed rectangular criterion on these samples does not constitute a separate continuous-uncertainty proof.

A distinct correlated-diagonal three-point reference yielded false-approval counts \(16,0,0,0\), and conservative-rejection counts \(0,0,3,3\), in the same method order. The latter three-point reference is explicitly \emph{finite-scenario}; it cannot certify an entire correlated interval. The analytical correlated interior counterexample in Proposition~\ref{prop:p13-counterexample} shows why a two-endpoint common-derating check may approve a case with an interior coordination failure. Conversely, a rectangular certificate is sufficient for a correlated subset of its box but may reject scenarios that are safe only because of their correlations.

\paragraph{Timing and interpretation.}
All implementations used the same stipulated time-domain relay integrator. The CI benchmark recorded approximately \(0.00040,0.00054,0.00056,0.00471\) seconds for the four complete \(90\)-scenario evaluations, respectively, on its particular hosted runner. These timings are illustrative, not machine-independent complexity measurements; repeated timing distributions and confidence intervals were not established. The benchmark is a controlled synthetic comparison, not an IEEE feeder switching or field-relay validation. It issues no physical restoration authorization.